\textbf{结论名称：}三个向量求和为零\(\iff\)首尾相连构成三角形。

\textbf{适用条件：}三个非零平面向量，两两不共线（或相应的有向线段可构成三角形）。

\textbf{变量说明：} \(\vec{a},\vec{b},\vec{c}\) 为平面内三个非零向量；用有向线段表示时，存在点 \(A,B,C\) 使 \(\vec{a}=\overrightarrow{AB},\vec{b}=\overrightarrow{BC},\vec{c}=\overrightarrow{CA}\)。

\textbf{核心公式：}
\[
\boxed{\vec{a} + \vec{b} + \vec{c} = \mathbf{0} \iff \text{以}\vec{a},\vec{b},\vec{c}\text{为边可首尾相连构成三角形}}
\]
或以有向线段形式：
\[
\boxed{\overrightarrow{AB} + \overrightarrow{BC} + \overrightarrow{CA} = \mathbf{0}}
\]

\textbf{使用场景：}将向量和的代数条件转化为几何图形（三角形），用于解向量模长、夹角问题；或由几何三角形反推向量线性关系。

\textbf{简单例子：}已知 \(|\vec{a}|=3,\ |\vec{b}|=4\)，且 \(\vec{a}+\vec{b}+\vec{c}=\mathbf{0}\)，三向量首尾相连构成直角三角形，求 \(|\vec{c}|\)。\\
解：由结论，\(\vec{a},\vec{b},\vec{c}\) 构成三角形。若 \(\vec{a}\perp\vec{b}\)，则由勾股定理得 \(|\vec{c}|=\sqrt{3^2+4^2}=5\)（向量 \(\vec{c}\) 为斜边）。